\section*{Prompt: Aggregate Constructs}
\begin{description}
\item[Pipeline step:] Harmonization --- Construct Grouping
\item[Models:] Gemini 2.5 Pro
\item[Source:] \texttt{misc/prompts.py} $\rightarrow$ \texttt{build\_aggregate\_grouping\_prompt('constructs', ...)}
\item[Called from:] \texttt{scripts/llm\_aggregator.py:group\_nodes\_llm()}
\end{description}
\begin{tcolorbox}[colback=yellow!8,colframe=yellow!50!black,left=4pt,right=4pt,top=4pt,bottom=4pt,arc=2pt]
\textbf{Note:} The \texttt{[... nodes injected here ...]} placeholder is replaced at runtime with a JSON array of construct nodes from all papers.
\end{tcolorbox}

\begin{tcolorbox}[breakable,enhanced,colback=blue!3,colframe=blue!40!black,title={Prompt},fonttitle=\bfseries,left=8pt,right=8pt,top=6pt,bottom=6pt]

Deduplicate constructs (latent LLM capabilities, risks, traits) by grouping rows that name the \textbf{exact same concept} under different wording.

Each input row has an \texttt{id}. Output only the grouping as \texttt{\{canonical\_name: [ids]\}}. All metadata is resolved from the IDs programmatically --- do NOT repeat definitions, papers, or snippets.

\textbf{Canonical name rules:} concise, community-standard noun phrase. Sentence case.
\begin{itemize}
\item \textbf{Drop trailing filler words} "ability", "capability", "capacity", "skill" when they are a separate word: "Reasoning ability" $\rightarrow$ "Reasoning", "Self-evolution capability" $\rightarrow$ "Self-evolution", "Retrieval capability" $\rightarrow$ "Retrieval"
\item \textbf{Keep compound words} where the suffix is integral: "Interpretability", "Editability", "Learnability", "Scalability" stay as-is
\end{itemize}

\textbf{Grouping rules:}
\begin{itemize}
\item Group the same concept said differently (e.g. "Reasoning" and "Reasoning Ability")
\item DO NOT group similar but distinct concepts (e.g. "Reasoning" vs "Logical Reasoning")
\item DO NOT group hierarchical variants (e.g. "Reasoning" vs "Mathematical Reasoning" --- keep separate)
\item DO NOT merge qualified/scoped variants into a bare parent (e.g. "Adversarial Robustness" $\neq$ "Robustness"; "In-context Learning" $\neq$ "Knowledge")
\item DO NOT merge distinct facets of a broader concept (e.g. "Mathematical reasoning", "Spatial reasoning", "Commonsense reasoning" are each their own group --- they are facets of "Reasoning", not synonyms)
\item Test: would a domain expert use the two names \textit{interchangeably}? If not, keep separate
\item When in doubt, keep separate
\end{itemize}

\textbf{Output format:}
\begin{verbatim}
{"groups": {"Reasoning": [0, 3], "Logical reasoning": [5]}, "ungrouped": [7, 12]}
\end{verbatim}
WRONG canonical names (trailing filler word) $\rightarrow$ RIGHT:
\begin{itemize}
\item "Self-evolution capability" $\rightarrow$ "Self-evolution"
\item "Generation capability" $\rightarrow$ "Generation"
\item "Retrieval capability" $\rightarrow$ "Retrieval"
\item "In-context learning ability" $\rightarrow$ "In-context learning"
\item "Implicit knowledge base capacity" $\rightarrow$ "Implicit knowledge base"
\end{itemize}

\texttt{ungrouped} lists IDs that don't match any other row. Every input ID must appear exactly once.

\textbf{Input rows:}
[... nodes injected here ...]

Return only the JSON object.
\end{tcolorbox}
